\documentclass[11pt,a4paper]{article}

\usepackage[utf8]{inputenc}
\usepackage[T1]{fontenc}
\usepackage[french]{babel}
\usepackage[a4paper,margin=2cm]{geometry}
\usepackage{amsmath,amssymb,amsthm,mathtools}
\usepackage{lmodern}
\usepackage{enumitem}
\usepackage{xcolor}
\usepackage{hyperref}
\usepackage{fancyhdr}
\usepackage{tikz}
\usepackage{booktabs}
\usepackage{array}

\hypersetup{
  colorlinks=true,
  linkcolor=black,
  urlcolor=blue
}

\setlength{\parindent}{0pt}
\setlength{\parskip}{0.5em}

\pagestyle{fancy}
\fancyhf{}
\fancyhead[L]{Poly de cours démontré}
\fancyhead[R]{Fonctions -- Lycée}
\fancyfoot[C]{\thepage}

\newtheorem{definition}{Définition}[section]
\newtheorem{propriete}[definition]{Propriété}
\newtheorem{theoreme}[definition]{Théorème}
\newtheorem{corollaire}[definition]{Corollaire}
\newtheorem*{exemple}{Exemple}
\newtheorem*{methode}{Méthode}
\newtheorem*{remarque}{Remarque}
\newtheorem*{admis}{Résultat admis}

\newcommand{\R}{\mathbb{R}}
\newcommand{\D}{\mathcal{D}}
\newcommand{\e}{\mathrm{e}}

\begin{document}

\begin{center}
{\LARGE \textbf{Poly de cours démontré -- Fonctions au lycée}}\\[0.5em]
{\large Seconde, Première, Terminale}\\[0.5em]
\end{center}

\tableofcontents
\newpage

\section{Notion de fonction}

\begin{definition}
Une \textbf{fonction} \(f\) définie sur un ensemble \(D \subset \R\) associe à tout réel \(x \in D\) un unique réel noté \(f(x)\).
\end{definition}

\begin{definition}
L'ensemble \(D\) s'appelle l'\textbf{ensemble de définition} de \(f\), noté \(\D_f\).
\end{definition}

\begin{definition}
Si \(f(a)=b\), alors :
\begin{itemize}
    \item \(b\) est l'\textbf{image} de \(a\),
    \item \(a\) est un \textbf{antécédent} de \(b\).
\end{itemize}
\end{definition}

\begin{propriete}
Une image est unique, mais un réel peut avoir zéro, un ou plusieurs antécédents.
\end{propriete}

\textbf{Démonstration.}  
Par définition d'une fonction, à un réel \(x\) de l'ensemble de définition est associé un \emph{unique} réel \(f(x)\), donc l'image est unique.

En revanche, plusieurs réels distincts peuvent avoir la même image. Par exemple, pour \(f(x)=x^2\), on a
\[
f(2)=4 \quad \text{et} \quad f(-2)=4.
\]
Le réel \(4\) a donc plusieurs antécédents. Il peut aussi n'en avoir aucun : par exemple, l'équation \(x^2=-1\) n'a pas de solution dans \(\R\). \(\square\)

\begin{exemple}
Pour \(f(x)=x^2\), l'image de \(3\) est \(9\), et les antécédents de \(9\) sont \(-3\) et \(3\).
\end{exemple}

\section{Définitions de base sur les fonctions}

\begin{definition}[Intervalle]
On appelle \textbf{intervalle} de \(\R\) un ensemble \(I\subset \R\) tel que :
\[
\forall a,b\in I,\ \forall x\in \R,\quad a\le x\le b \Rightarrow x\in I.
\]
Autrement dit, dès que deux réels appartiennent à \(I\), alors tous les réels compris entre eux appartiennent aussi à \(I\).
\end{definition}

\begin{exemple}
Les ensembles
\[
[0,1],\qquad ]-\infty,3],\qquad \R
\]
sont des intervalles.
\end{exemple}

\begin{definition}[Ensemble majoré, minoré, borné]
Soit \(A\subset \R\).
\begin{itemize}
    \item On dit que \(A\) est \textbf{majoré} s'il existe un réel \(M\) tel que
    \[
    \forall x\in A,\quad x\le M.
    \]
    \item On dit que \(A\) est \textbf{minoré} s'il existe un réel \(m\) tel que
    \[
    \forall x\in A,\quad x\ge m.
    \]
    \item On dit que \(A\) est \textbf{borné} s'il est à la fois majoré et minoré.
\end{itemize}
\end{definition}

\begin{exemple}
L'intervalle \([2,5]\) est borné.  
L'intervalle \([0,+\infty[\) est minoré mais n'est pas majoré.
\end{exemple}

\begin{definition}[Fonction croissante]
Soit \(f\) une fonction définie sur un intervalle \(I\).  
On dit que \(f\) est \textbf{croissante} sur \(I\) si
\[
\forall x_1,x_2\in I,\quad x_1\le x_2 \Rightarrow f(x_1)\le f(x_2).
\]
\end{definition}

\begin{definition}[Fonction strictement croissante]
Soit \(f\) une fonction définie sur un intervalle \(I\).  
On dit que \(f\) est \textbf{strictement croissante} sur \(I\) si
\[
\forall x_1,x_2\in I,\quad x_1<x_2 \Rightarrow f(x_1)<f(x_2).
\]
\end{definition}

\begin{definition}[Fonction décroissante]
Soit \(f\) une fonction définie sur un intervalle \(I\).  
On dit que \(f\) est \textbf{décroissante} sur \(I\) si
\[
\forall x_1,x_2\in I,\quad x_1\le x_2 \Rightarrow f(x_1)\ge f(x_2).
\]
\end{definition}

\begin{definition}[Fonction strictement décroissante]
Soit \(f\) une fonction définie sur un intervalle \(I\).  
On dit que \(f\) est \textbf{strictement décroissante} sur \(I\) si
\[
\forall x_1,x_2\in I,\quad x_1<x_2 \Rightarrow f(x_1)>f(x_2).
\]
\end{definition}

\begin{definition}[Fonction monotone]
Une fonction est dite \textbf{monotone} sur un intervalle lorsqu'elle y est soit croissante, soit décroissante.
\end{definition}

\begin{remarque}
Pour montrer qu'une fonction est croissante, on prend deux réels \(x_1\) et \(x_2\) tels que \(x_1\le x_2\), puis on compare \(f(x_1)\) et \(f(x_2)\).

Pour montrer qu'une fonction est strictement croissante, on prend deux réels \(x_1<x_2\), puis on montre que \(f(x_1)<f(x_2)\).
\end{remarque}

\begin{exemple}
Soit \(f(x)=3x+1\).  
Montrons que \(f\) est strictement croissante sur \(\R\).

Soient \(x_1<x_2\). Alors
\[
f(x_2)-f(x_1)=(3x_2+1)-(3x_1+1)=3(x_2-x_1).
\]
Or \(x_2-x_1>0\), donc
\[
3(x_2-x_1)>0.
\]
Ainsi
\[
f(x_2)-f(x_1)>0,
\]
donc
\[
f(x_2)>f(x_1).
\]
La fonction \(f\) est donc strictement croissante sur \(\R\). \(\square\)
\end{exemple}

\begin{definition}[Maximum]
Soit \(f\) définie sur un ensemble \(A\subset \R\).  
On dit que \(f\) admet un \textbf{maximum} en \(a\in A\) si
\[
\forall x\in A,\quad f(x)\le f(a).
\]
Le réel \(f(a)\) est alors appelé le maximum de \(f\) sur \(A\).
\end{definition}

\begin{definition}[Minimum]
Soit \(f\) définie sur un ensemble \(A\subset \R\).  
On dit que \(f\) admet un \textbf{minimum} en \(a\in A\) si
\[
\forall x\in A,\quad f(x)\ge f(a).
\]
Le réel \(f(a)\) est alors appelé le minimum de \(f\) sur \(A\).
\end{definition}

\begin{definition}[Extremum]
Un \textbf{extremum} est soit un maximum, soit un minimum.
\end{definition}

\begin{exemple}
Pour \(f(x)=x^2\) sur \(\R\), on a
\[
\forall x\in\R,\quad x^2\ge 0=f(0).
\]
Donc \(f\) admet un minimum en \(0\), égal à \(0\).
\end{exemple}

\section{Ensemble de définition}

\begin{propriete}
Pour déterminer l'ensemble de définition d'une fonction, on impose que toutes les expressions qui interviennent aient un sens.
\end{propriete}

\subsection{Cas usuels}

\begin{itemize}
    \item Un polynôme est défini sur \(\R\).
    \item Un quotient \(\dfrac{u(x)}{v(x)}\) est défini lorsque \(v(x)\neq 0\).
    \item Une racine \(\sqrt{u(x)}\) est définie lorsque \(u(x)\ge 0\).
    \item Un logarithme \(\ln(u(x))\) est défini lorsque \(u(x)>0\).
\end{itemize}

\begin{exemple}
\[
f(x)=\frac{1}{x-2}
\qquad \Rightarrow \qquad \D_f=\R\setminus\{2\}.
\]
\end{exemple}

\begin{exemple}
\[
g(x)=\sqrt{3x+1}
\]
est définie si \(3x+1\ge 0\), soit
\[
x\ge -\frac13.
\]
Donc
\[
\D_g=\left[-\frac13,+\infty\right[.
\]
\end{exemple}

\section{Fonctions usuelles}

\subsection{Fonction affine}

\begin{definition}
Une fonction affine est une fonction de la forme
\[
f(x)=ax+b
\]
où \(a,b\in\R\).
\end{definition}

\begin{propriete}
\begin{itemize}
    \item Si \(a>0\), \(f\) est strictement croissante sur \(\R\).
    \item Si \(a<0\), \(f\) est strictement décroissante sur \(\R\).
    \item Si \(a=0\), \(f\) est constante.
\end{itemize}
\end{propriete}

\textbf{Démonstration.}  
Soient \(x_1<x_2\). Alors
\[
f(x_2)-f(x_1)=(ax_2+b)-(ax_1+b)=a(x_2-x_1).
\]
Or \(x_2-x_1>0\).

\begin{itemize}
    \item Si \(a>0\), alors \(a(x_2-x_1)>0\), donc \(f(x_2)>f(x_1)\) : \(f\) est strictement croissante.
    \item Si \(a<0\), alors \(a(x_2-x_1)<0\), donc \(f(x_2)<f(x_1)\) : \(f\) est strictement décroissante.
    \item Si \(a=0\), alors \(f(x)=b\) pour tout \(x\), donc \(f\) est constante.
\end{itemize}
\(\square\)

\subsection{Fonction carré}

\[
f(x)=x^2
\]

\begin{propriete}
La fonction carré est décroissante sur \(]-\infty,0]\) et croissante sur \([0,+\infty[\).
\end{propriete}

\textbf{Démonstration.}  
Soient \(x_1<x_2\).

\emph{Sur \([0,+\infty[\).}  
On a \(x_1\ge 0\) et \(x_2\ge 0\). Alors
\[
x_2^2-x_1^2=(x_2-x_1)(x_2+x_1).
\]
Or \(x_2-x_1>0\) et \(x_2+x_1\ge 0\), donc
\[
x_2^2-x_1^2\ge 0.
\]
Ainsi \(x_1^2\le x_2^2\), donc la fonction est croissante sur \([0,+\infty[\).

\emph{Sur \(]-\infty,0]\).}  
Si \(x_1<x_2\le 0\), alors \(-x_1>-x_2\ge 0\). Comme la fonction carré prend la même valeur en \(x\) et en \(-x\), et qu'elle est croissante sur \([0,+\infty[\), on obtient
\[
(-x_1)^2\ge (-x_2)^2,
\]
donc
\[
x_1^2\ge x_2^2.
\]
La fonction est donc décroissante sur \(]-\infty,0]\). \(\square\)

\subsection{Fonction cube}

\[
f(x)=x^3
\]

\begin{propriete}
La fonction cube est strictement croissante sur \(\R\).
\end{propriete}

\textbf{Démonstration.}  
Soient \(x_1<x_2\). Alors
\[
x_2^3-x_1^3=(x_2-x_1)(x_2^2+x_1x_2+x_1^2).
\]
On a \(x_2-x_1>0\). De plus
\[
x_2^2+x_1x_2+x_1^2=\left(x_2+\frac{x_1}{2}\right)^2+\frac{3}{4}x_1^2\ge 0.
\]
Cette quantité n'est nulle que si \(x_1=x_2=0\), ce qui est impossible puisque \(x_1<x_2\). Donc elle est strictement positive. Ainsi
\[
x_2^3-x_1^3>0,
\]
donc \(x_1^3<x_2^3\). \(\square\)

\subsection{Fonction inverse}

\[
f(x)=\frac{1}{x}
\]

\begin{propriete}
La fonction inverse est définie sur \(\R\setminus\{0\}\) et strictement décroissante sur \(]-\infty,0[\) et sur \(]0,+\infty[\).
\end{propriete}

\textbf{Démonstration.}  
Soient \(x_1<x_2\), avec \(x_1\) et \(x_2\) de même signe.

On calcule
\[
\frac{1}{x_2}-\frac{1}{x_1}=\frac{x_1-x_2}{x_1x_2}.
\]

\begin{itemize}
    \item Si \(0<x_1<x_2\), alors \(x_1-x_2<0\) et \(x_1x_2>0\), donc
    \[
    \frac{1}{x_2}-\frac{1}{x_1}<0.
    \]
    Ainsi \(\frac1{x_2}<\frac1{x_1}\), donc la fonction est décroissante sur \(]0,+\infty[\).

    \item Si \(x_1<x_2<0\), alors encore \(x_1-x_2<0\) et \(x_1x_2>0\), donc
    \[
    \frac{1}{x_2}-\frac{1}{x_1}<0.
    \]
    La fonction est aussi décroissante sur \(]-\infty,0[\).
\end{itemize}
\(\square\)

\subsection{Fonction racine carrée}

\[
f(x)=\sqrt{x}
\]

\begin{propriete}
La fonction racine carrée est définie sur \([0,+\infty[\) et strictement croissante sur cet intervalle.
\end{propriete}

\textbf{Démonstration.}  
Soient \(0\le x_1<x_2\). Supposons par l'absurde que \(\sqrt{x_1}\ge \sqrt{x_2}\).  
Comme les deux membres sont positifs, on peut élever au carré et obtenir
\[
x_1\ge x_2,
\]
ce qui contredit \(x_1<x_2\). Donc
\[
\sqrt{x_1}<\sqrt{x_2}.
\]
La fonction est strictement croissante. \(\square\)

\section{Parité}

\begin{definition}
Une fonction \(f\) définie sur un ensemble symétrique par rapport à \(0\) est :
\begin{itemize}
    \item \textbf{paire} si \(\forall x,\ f(-x)=f(x)\),
    \item \textbf{impaire} si \(\forall x,\ f(-x)=-f(x)\).
\end{itemize}
\end{definition}

\begin{propriete}
\begin{itemize}
    \item La courbe d'une fonction paire est symétrique par rapport à l'axe des ordonnées.
    \item La courbe d'une fonction impaire est symétrique par rapport à l'origine.
\end{itemize}
\end{propriete}

\textbf{Démonstration.}  
Soit un point \(M(x,f(x))\) de la courbe.

\begin{itemize}
    \item Si \(f\) est paire, alors \(f(-x)=f(x)\). Le point \((-x,f(x))\) appartient aussi à la courbe. C'est le symétrique de \(M\) par rapport à l'axe des ordonnées.
    \item Si \(f\) est impaire, alors \(f(-x)=-f(x)\). Le point \((-x,-f(x))\) appartient à la courbe. C'est le symétrique de \(M\) par rapport à l'origine.
\end{itemize}
\(\square\)

\section{Limites}

\begin{definition}
La limite d'une fonction décrit son comportement lorsque \(x\) s'approche d'une valeur donnée ou devient très grand en valeur absolue.
\end{definition}

\subsection{Limites usuelles}

\[
\lim_{x\to +\infty}x=+\infty,\qquad
\lim_{x\to +\infty}x^2=+\infty,\qquad
\lim_{x\to +\infty}\frac1x=0.
\]

\[
\lim_{x\to 0^+}\frac1x=+\infty,\qquad
\lim_{x\to 0^-}\frac1x=-\infty.
\]

\begin{propriete}
\[
\lim_{x\to +\infty}\frac1x=0.
\]
\end{propriete}

\textbf{Justification.}  
Quand \(x\) devient de plus en plus grand, le dénominateur de \(\frac1x\) devient immense, donc la fraction devient de plus en plus petite en valeur absolue. On peut aussi écrire : pour tout \(\varepsilon>0\), si \(x>\frac1\varepsilon\), alors
\[
\left|\frac1x\right|<\varepsilon.
\]
Donc \(\frac1x\to 0\). \(\square\)

\subsection{Asymptotes}

\begin{definition}
La droite \(x=a\) est une \textbf{asymptote verticale} si
\[
\lim_{x\to a}f(x)=+\infty \quad \text{ou} \quad -\infty.
\]
\end{definition}

\begin{definition}
La droite \(y=\ell\) est une \textbf{asymptote horizontale} si
\[
\lim_{x\to \pm\infty}f(x)=\ell.
\]
\end{definition}

\begin{exemple}
Pour \(f(x)=\dfrac1x\), la droite \(x=0\) est asymptote verticale et la droite \(y=0\) est asymptote horizontale.
\end{exemple}

\section{Continuité}

\begin{definition}
Une fonction \(f\) est \textbf{continue} en \(a\) si
\[
\lim_{x\to a}f(x)=f(a).
\]
\end{definition}

\begin{propriete}
Les polynômes sont continus sur \(\R\). Les fonctions rationnelles sont continues sur leur ensemble de définition.
\end{propriete}

\begin{remarque}
Au lycée, cette propriété est en général admise ou déjà connue pour les fonctions usuelles.
\end{remarque}

\subsection{Théorème des valeurs intermédiaires}

\begin{theoreme}[TVI]
Si \(f\) est continue sur \([a,b]\), alors pour tout réel \(k\) compris entre \(f(a)\) et \(f(b)\), il existe au moins un réel \(c\in[a,b]\) tel que
\[
f(c)=k.
\]
\end{theoreme}

\begin{admis}
Le TVI est un théorème fondamental du cours, mais sa démonstration rigoureuse n'est généralement pas exigée au lycée.
\end{admis}

\begin{corollaire}
Si \(f\) est continue sur \([a,b]\) et si \(f(a)\) et \(f(b)\) sont de signes contraires, alors l'équation \(f(x)=0\) admet au moins une solution dans \([a,b]\).
\end{corollaire}

\textbf{Démonstration.}  
Si \(f(a)\) et \(f(b)\) sont de signes contraires, alors \(0\) est compris entre \(f(a)\) et \(f(b)\).  
Par le théorème des valeurs intermédiaires, il existe donc \(c\in[a,b]\) tel que
\[
f(c)=0.
\]
\(\square\)

\begin{exemple}
Soit
\[
f(x)=x^3-x-1.
\]
C'est un polynôme, donc \(f\) est continue sur \([1,2]\). De plus
\[
f(1)=-1,\qquad f(2)=5.
\]
Comme \(f(1)\) et \(f(2)\) sont de signes contraires, l'équation \(f(x)=0\) admet au moins une solution dans \([1,2]\).
\end{exemple}

\section{Dérivation}

\begin{definition}
Soit \(f\) définie sur un intervalle \(I\) et \(a\in I\). On dit que \(f\) est \textbf{dérivable} en \(a\) si la limite
\[
\lim_{h\to 0}\frac{f(a+h)-f(a)}{h}
\]
existe et est finie. Cette limite est notée \(f'(a)\).
\end{definition}

\begin{definition}
La dérivée \(f'(a)\) est le coefficient directeur de la tangente à la courbe de \(f\) au point d'abscisse \(a\).
\end{definition}

\subsection{Tangente}

\begin{propriete}
Si \(f\) est dérivable en \(a\), une équation de la tangente en \(a\) est
\[
y=f(a)+f'(a)(x-a).
\]
\end{propriete}

\textbf{Justification.}  
Une droite de coefficient directeur \(m\) passant par \((a,f(a))\) a pour équation
\[
y-f(a)=m(x-a).
\]
Or ici \(m=f'(a)\). D'où
\[
y=f(a)+f'(a)(x-a).
\]
\(\square\)

\subsection{Formules de dérivation}

\begin{propriete}
\[
(x^2)'=2x,\qquad (x^3)'=3x^2,\qquad \left(\frac1x\right)'=-\frac1{x^2}.
\]
\end{propriete}

\textbf{Démonstration de \((x^2)'=2x\).}  
Soit \(f(x)=x^2\). Alors
\[
\frac{f(a+h)-f(a)}{h}=\frac{(a+h)^2-a^2}{h}
=\frac{a^2+2ah+h^2-a^2}{h}
=\frac{2ah+h^2}{h}=2a+h.
\]
Quand \(h\to 0\), on obtient
\[
f'(a)=2a.
\]
Donc \((x^2)'=2x\). \(\square\)

\textbf{Démonstration de \((x^3)'=3x^2\).}  
Soit \(f(x)=x^3\). Alors
\[
\frac{f(a+h)-f(a)}{h}
=\frac{(a+h)^3-a^3}{h}
=\frac{a^3+3a^2h+3ah^2+h^3-a^3}{h}
=3a^2+3ah+h^2.
\]
Quand \(h\to 0\), on obtient
\[
f'(a)=3a^2.
\]
Donc \((x^3)'=3x^2\). \(\square\)

\textbf{Démonstration de \(\left(\frac1x\right)'=-\frac1{x^2}\).}  
Soit \(f(x)=\frac1x\), avec \(a\neq 0\). Alors
\[
\frac{f(a+h)-f(a)}{h}
=
\frac{\frac1{a+h}-\frac1a}{h}
=
\frac{\frac{a-(a+h)}{a(a+h)}}{h}
=
\frac{-h}{a(a+h)h}
=
-\frac1{a(a+h)}.
\]
Quand \(h\to 0\), on obtient
\[
f'(a)=-\frac1{a^2}.
\]
Donc
\[
\left(\frac1x\right)'=-\frac1{x^2}.
\]
\(\square\)

\subsection{Somme et produit}

\begin{propriete}
Si \(u\) et \(v\) sont dérivables, alors
\[
(u+v)'=u'+v'
\]
et
\[
(uv)'=u'v+uv'.
\]
\end{propriete}

\textbf{Démonstration de la somme.}  
On écrit
\[
\frac{(u+v)(a+h)-(u+v)(a)}{h}
=
\frac{u(a+h)-u(a)}{h}+\frac{v(a+h)-v(a)}{h}.
\]
En passant à la limite quand \(h\to 0\), on obtient
\[
(u+v)'(a)=u'(a)+v'(a).
\]
\(\square\)

\textbf{Démonstration du produit.}  
On part de
\[
\frac{u(a+h)v(a+h)-u(a)v(a)}{h}.
\]
On ajoute et retranche \(u(a+h)v(a)\) :
\[
\frac{u(a+h)v(a+h)-u(a+h)v(a)+u(a+h)v(a)-u(a)v(a)}{h}.
\]
On factorise :
\[
u(a+h)\frac{v(a+h)-v(a)}{h}+v(a)\frac{u(a+h)-u(a)}{h}.
\]
Quand \(h\to 0\), \(u(a+h)\to u(a)\) car une fonction dérivable est continue, donc
\[
(uv)'(a)=u(a)v'(a)+v(a)u'(a).
\]
\(\square\)

\section{Lien entre dérivée et variations}

\begin{theoreme}
Soit \(f\) dérivable sur un intervalle \(I\).
\begin{itemize}
    \item Si \(f'(x)\ge 0\) sur \(I\), alors \(f\) est croissante sur \(I\).
    \item Si \(f'(x)\le 0\) sur \(I\), alors \(f\) est décroissante sur \(I\).
\end{itemize}
\end{theoreme}

\begin{demonstration}
Soit \(f\) une fonction dérivable en un réel \(a\).

\begin{itemize}
    \item Si \(f'(a)>0\), alors il existe un voisinage de \(a\) sur lequel :
    \[
    x<a \Rightarrow f(x)<f(a)
    \qquad \text{et} \qquad
    x>a \Rightarrow f(x)>f(a).
    \]
    \item Si \(f'(a)<0\), alors il existe un voisinage de \(a\) sur lequel :
    \[
    x<a \Rightarrow f(x)>f(a)
    \qquad \text{et} \qquad
    x>a \Rightarrow f(x)<f(a).
    \]
\end{itemize}
\end{theoreme}

\begin{proof}
Supposons \(f'(a)>0\).

Par définition de la dérivée en \(a\),
\[
f'(a)=\lim_{x\to a}\frac{f(x)-f(a)}{x-a}.
\]
Posons
\[
L=f'(a).
\]
On a donc \(L>0\).

Comme
\[
\lim_{x\to a}\frac{f(x)-f(a)}{x-a}=L,
\]
on peut appliquer la définition de la limite avec
\[
\varepsilon=\frac{L}{2}>0.
\]
Il existe alors \(\delta>0\) tel que, pour tout \(x\) vérifiant
\[
0<|x-a|<\delta,
\]
on ait
\[
\left|\frac{f(x)-f(a)}{x-a}-L\right|<\frac{L}{2}.
\]
Cette inégalité entraîne
\[
-\frac{L}{2}
<
\frac{f(x)-f(a)}{x-a}-L
<
\frac{L}{2}.
\]
En ajoutant \(L\) dans chaque membre, on obtient
\[
\frac{L}{2}
<
\frac{f(x)-f(a)}{x-a}
<
\frac{3L}{2}.
\]
En particulier,
\[
\frac{f(x)-f(a)}{x-a}>0
\qquad \text{dès que } 0<|x-a|<\delta.
\]

Ainsi, pour tout \(x\) suffisamment proche de \(a\), le quotient
\[
\frac{f(x)-f(a)}{x-a}
\]
est strictement positif. Le numérateur \(f(x)-f(a)\) a donc le même signe que le dénominateur \(x-a\).

\begin{itemize}
    \item Si \(a-\delta<x<a\), alors \(x-a<0\). Comme
    \[
    \frac{f(x)-f(a)}{x-a}>0,
    \]
    on doit avoir
    \[
    f(x)-f(a)<0,
    \]
    donc
    \[
    f(x)<f(a).
    \]

    \item Si \(a<x<a+\delta\), alors \(x-a>0\). Comme
    \[
    \frac{f(x)-f(a)}{x-a}>0,
    \]
    on obtient
    \[
    f(x)-f(a)>0,
    \]
    donc
    \[
    f(x)>f(a).
    \]
\end{itemize}

Cela prouve le premier point.

\vspace{0.5em}

Le cas \(f'(a)<0\) se traite de la même manière.  
En effet, si
\[
\lim_{x\to a}\frac{f(x)-f(a)}{x-a}=f'(a)<0,
\]
alors le quotient
\[
\frac{f(x)-f(a)}{x-a}
\]
est strictement négatif au voisinage de \(a\). Le numérateur \(f(x)-f(a)\) est donc de signe opposé à \(x-a\).

On en déduit :
\begin{itemize}
    \item si \(x<a\), alors \(x-a<0\), donc \(f(x)-f(a)>0\), d'où
    \[
    f(x)>f(a);
    \]
    \item si \(x>a\), alors \(x-a>0\), donc \(f(x)-f(a)<0\), d'où
    \[
    f(x)<f(a).
    \]
\end{itemize}
\end{demonstration}

\begin{corollaire}
Si \(f'(x)>0\) sur \(I\), alors \(f\) est strictement croissante sur \(I\).  
Si \(f'(x)<0\) sur \(I\), alors \(f\) est strictement décroissante sur \(I\).
\end{corollaire}

\begin{exemple}
Soit
\[
f(x)=x^2-4x+1.
\]
Alors
\[
f'(x)=2x-4=2(x-2).
\]
Donc :
\begin{itemize}
    \item \(f'(x)<0\) si \(x<2\),
    \item \(f'(x)=0\) si \(x=2\),
    \item \(f'(x)>0\) si \(x>2\).
\end{itemize}
Ainsi \(f\) est décroissante sur \(]-\infty,2]\) et croissante sur \([2,+\infty[\).
\end{exemple}

\section{Extremums}

\begin{definition}
Soit \(f\) définie sur un ensemble \(I\). On dit que \(f(a)\) est :
\begin{itemize}
    \item un \textbf{maximum} sur \(I\) si \(\forall x\in I,\ f(x)\le f(a)\),
    \item un \textbf{minimum} sur \(I\) si \(\forall x\in I,\ f(x)\ge f(a)\).
\end{itemize}
\end{definition}

\begin{propriete}
Lorsqu'une fonction est décroissante puis croissante, elle admet un minimum au point de changement. Lorsqu'elle est croissante puis décroissante, elle admet un maximum.
\end{propriete}

\textbf{Justification.}  
Si \(f\) décroît jusqu'en \(a\), alors pour \(x\le a\), on a \(f(x)\ge f(a)\).  
Si ensuite \(f\) croît à partir de \(a\), alors pour \(x\ge a\), on a aussi \(f(x)\ge f(a)\).  
Donc pour tout \(x\), \(f(x)\ge f(a)\), ce qui montre que \(f(a)\) est un minimum. \(\square\)

\section{Convexité}

\begin{definition}
Une fonction est \textbf{convexe} lorsque sa courbe est tournée vers le haut.  
Une fonction est \textbf{concave} lorsque sa courbe est tournée vers le bas.
\end{definition}

\begin{propriete}
Si \(f\) est deux fois dérivable :
\begin{itemize}
    \item si \(f''(x)\ge 0\), alors \(f\) est convexe ;
    \item si \(f''(x)\le 0\), alors \(f\) est concave.
\end{itemize}
\end{propriete}

\begin{admis}
Ce critère est utilisé au lycée, mais sa démonstration complète n'est généralement pas exigée.
\end{admis}

\begin{definition}
Un \textbf{point d'inflexion} est un point où la courbe change de convexité.
\end{definition}

\begin{propriete}
Si \(f''\) change de signe en \(a\), alors la courbe de \(f\) admet un point d'inflexion d'abscisse \(a\).
\end{propriete}

\textbf{Justification.}  
Si \(f''\) est négative d'un côté de \(a\) puis positive de l'autre, la fonction passe d'un comportement concave à un comportement convexe. Il y a donc changement de convexité en \(a\), ce qui correspond à un point d'inflexion. \(\square\)

\begin{exemple}
Soit
\[
f(x)=x^3.
\]
On a
\[
f'(x)=3x^2,\qquad f''(x)=6x.
\]
Donc \(f''(x)<0\) pour \(x<0\) et \(f''(x)>0\) pour \(x>0\). Il y a changement de signe en \(0\), donc \(0\) est abscisse d'un point d'inflexion.
\end{exemple}

\section{Fonction exponentielle}

\begin{definition}
La fonction exponentielle est la fonction \(x\mapsto \e^x\).
\end{definition}

\begin{propriete}
Pour tous réels \(a,b\),
\[
\e^{a+b}=\e^a\e^b,\qquad \e^0=1,\qquad \e^{-a}=\frac1{\e^a}.
\]
\end{propriete}

\begin{admis}
Ces propriétés caractérisent l'exponentielle et sont admises dans le cadre du lycée.
\end{admis}

\begin{propriete}
\[
(\e^x)'=\e^x.
\]
\end{propriete}

\begin{admis}
La dérivée de l'exponentielle est un résultat fondamental du cours, admis au lycée.
\end{admis}

\begin{corollaire}
La fonction exponentielle est strictement croissante sur \(\R\).
\end{corollaire}

\textbf{Démonstration.}  
Pour tout réel \(x\),
\[
(\e^x)'=\e^x.
\]
Or \(\e^x>0\) pour tout \(x\). Donc la dérivée est strictement positive sur \(\R\). La fonction exponentielle est donc strictement croissante. \(\square\)

\begin{propriete}
\[
\lim_{x\to +\infty}\e^x=+\infty,\qquad \lim_{x\to -\infty}\e^x=0.
\]
\end{propriete}

\begin{remarque}
Ces limites font partie des résultats usuels à connaître.
\end{remarque}

\section{Fonction logarithme népérien}

\begin{definition}
La fonction logarithme népérien, notée \(\ln\), est définie sur \(]0,+\infty[\). Elle est la fonction réciproque de l'exponentielle.
\end{definition}

\begin{propriete}
Pour tous \(a,b>0\),
\[
\ln(ab)=\ln(a)+\ln(b),\qquad
\ln\left(\frac{a}{b}\right)=\ln(a)-\ln(b).
\]
\end{propriete}

\begin{admis}
Ces propriétés sont admises au lycée.
\end{admis}

\begin{propriete}
\[
(\ln x)'=\frac1x \quad \text{sur } ]0,+\infty[.
\]
\end{propriete}

\begin{admis}
La dérivée du logarithme est un résultat fondamental admis dans le cadre du cours.
\end{admis}

\begin{corollaire}
La fonction \(\ln\) est strictement croissante sur \(]0,+\infty[\).
\end{corollaire}

\textbf{Démonstration.}  
Pour \(x>0\), on a
\[
(\ln x)'=\frac1x>0.
\]
Donc la fonction \(\ln\) est strictement croissante sur \(]0,+\infty[\). \(\square\)

\begin{propriete}
\[
\lim_{x\to 0^+}\ln x=-\infty,\qquad \lim_{x\to +\infty}\ln x=+\infty.
\]
\end{propriete}

\section{Fonctions composées}

\begin{definition}
Si \(u\) et \(v\) sont deux fonctions telles que l'image de \(u\) est contenue dans l'ensemble de définition de \(v\), alors
\[
(v\circ u)(x)=v(u(x)).
\]
\end{definition}

\begin{propriete}
Si \(u\) et \(v\) sont dérivables, alors
\[
(v\circ u)'(x)=v'(u(x))\cdot u'(x).
\]
\end{propriete}

\begin{admis}
La formule de dérivation des fonctions composées est utilisée au lycée, mais sa démonstration complète n'est pas toujours exigée.
\end{admis}

\begin{exemple}
Pour
\[
f(x)=\ln(2x+3),
\]
on a
\[
f'(x)=\frac{2}{2x+3}.
\]
\end{exemple}

\section{Suites définies par une fonction}

\begin{definition}
Une suite peut être définie par récurrence par une relation du type
\[
u_{n+1}=f(u_n).
\]
\end{definition}

\begin{exemple}
Soit
\[
u_{n+1}=\frac{u_n+2}{2},\qquad u_0=6.
\]
Si la suite converge vers un réel \(\ell\), alors en passant à la limite dans la relation
\[
u_{n+1}=\frac{u_n+2}{2},
\]
on obtient
\[
\ell=\frac{\ell+2}{2}.
\]
Donc
\[
2\ell=\ell+2 \quad \Rightarrow \quad \ell=2.
\]
Ainsi, si la suite converge, sa limite ne peut être que \(2\).
\end{exemple}

\section{Comparaison des croissances}

\begin{propriete}
Quand \(x\to +\infty\),
\[
\ln x \ll x^\alpha \ll \e^x \qquad (\alpha>0).
\]
\end{propriete}

\begin{admis}
Cette comparaison des croissances est connue au lycée, mais sa démonstration générale ne fait pas partie des démonstrations exigibles.
\end{admis}

\begin{corollaire}
\[
\lim_{x\to +\infty}\frac{\ln x}{x}=0,
\qquad
\lim_{x\to +\infty}\frac{x^2}{\e^x}=0.
\]
\end{corollaire}

\textbf{Justification.}  
Ces limites traduisent exactement que le logarithme croît beaucoup moins vite qu'une puissance, et qu'une puissance croît beaucoup moins vite que l'exponentielle. \(\square\)

\section{Méthode générale d'étude d'une fonction}

\begin{methode}
Pour étudier une fonction, on suit souvent les étapes suivantes :
\begin{enumerate}[label=\arabic*.]
    \item déterminer l'ensemble de définition ;
    \item calculer les limites aux bornes du domaine ;
    \item calculer la dérivée ;
    \item étudier le signe de la dérivée ;
    \item en déduire les variations ;
    \item déterminer les extremums ;
    \item étudier éventuellement la convexité avec \(f''\).
\end{enumerate}
\end{methode}

\section{Exemple complet d'étude}

\[
f(x)=x^2-4x+1
\]

\textbf{1. Domaine.}  
La fonction est un polynôme, donc elle est définie sur \(\R\).

\textbf{2. Dérivée.}
\[
f'(x)=2x-4=2(x-2).
\]

\textbf{3. Signe de la dérivée.}
\[
f'(x)<0 \text{ si } x<2,\qquad f'(2)=0,\qquad f'(x)>0 \text{ si } x>2.
\]

\textbf{4. Variations.}  
Donc \(f\) décroît sur \(]-\infty,2]\) puis croît sur \([2,+\infty[\).

\textbf{5. Minimum.}
\[
f(2)=4-8+1=-3.
\]
La fonction admet donc un minimum égal à \(-3\), atteint en \(x=2\).

\textbf{6. Convexité.}
\[
f''(x)=2>0.
\]
Donc \(f\) est convexe sur \(\R\).

\section{Exercices rapides}

\subsection*{Exercice 1}
Déterminer l'ensemble de définition de
\[
f(x)=\frac{\sqrt{x+1}}{x-3}.
\]

\subsection*{Exercice 2}
Montrer que la fonction \(x\mapsto 5x-2\) est strictement croissante.

\subsection*{Exercice 3}
Calculer la dérivée de
\[
f(x)=x^3+2x^2-1.
\]

\subsection*{Exercice 4}
Étudier les variations de
\[
f(x)=x^2-6x+5.
\]

\subsection*{Exercice 5}
Étudier la convexité de
\[
f(x)=x^3-3x.
\]

\section*{Réponses rapides}

\begin{itemize}
    \item Exercice 1 :
    \[
    \D_f=[-1,+\infty[\setminus\{3\}.
    \]

    \item Exercice 2 :
    si \(x_1<x_2\), alors
    \[
    (5x_2-2)-(5x_1-2)=5(x_2-x_1)>0,
    \]
    donc la fonction est strictement croissante.

    \item Exercice 3 :
    \[
    f'(x)=3x^2+4x.
    \]

    \item Exercice 4 :
    \[
    f'(x)=2x-6=2(x-3).
    \]
    Donc \(f\) décroît sur \(]-\infty,3]\), croît sur \([3,+\infty[\), avec minimum
    \[
    f(3)=9-18+5=-4.
    \]

    \item Exercice 5 :
    \[
    f''(x)=6x.
    \]
    Donc \(f\) est concave sur \(]-\infty,0[\), convexe sur \(]0,+\infty[\), avec point d'inflexion en \(0\).
\end{itemize}

\section*{Bilan}

Dans ce poly :
\begin{itemize}
    \item les résultats élémentaires sur les fonctions usuelles, la dérivation, les variations et les extremums sont démontrés ;
    \item les grands théorèmes de structure du lycée (TVI, critères avancés de convexité, comparaison générale des croissances) sont clairement signalés comme admis ;
    \item les méthodes essentielles d'étude sont illustrées par des exemples.
\end{itemize}

\end{document}